\begin{tikzpicture}[x=1cm,y=1cm,line width=.6pt,font=\normalsize,>=Stealth]
% hexagonal nut: regular hexagon plan (side R=2.2), oblique depth factor 0.45,
% prism height H=1.2; through-hole radius 1.1 (diameter 20 vs side 20, height 10)
\coordinate (T0) at (2.2,1.6);
\coordinate (T60) at (1.1,2.458);
\coordinate (T120) at (-1.1,2.458);
\coordinate (T180) at (-2.2,1.6);
\coordinate (T240) at (-1.1,0.742);
\coordinate (T300) at (1.1,0.742);
\coordinate (B0) at (2.2,0.4);
\coordinate (B60) at (1.1,1.258);
\coordinate (B120) at (-1.1,1.258);
\coordinate (B180) at (-2.2,0.4);
\coordinate (B240) at (-1.1,-0.458);
\coordinate (B300) at (1.1,-0.458);
% top face
\draw[solid] (T0)--(T60)--(T120)--(T180)--(T240)--(T300)--cycle;
% visible verticals and front bottom edges
\draw[solid] (T180)--(B180);
\draw[solid] (T240)--(B240);
\draw[solid] (T300)--(B300);
\draw[solid] (T0)--(B0);
\draw[solid] (B180)--(B240)--(B300)--(B0);
% hidden verticals and back bottom edges
\draw[solid,dashed] (T60)--(B60);
\draw[solid,dashed] (T120)--(B120);
\draw[solid,dashed] (B0)--(B60)--(B120)--(B180);
% through-hole: top rim fully visible, walls and bottom rim hidden
\draw[solid] (0,1.6) ellipse (1.1 and 0.495);
\draw[solid,dashed] (1.1,1.6)--(1.1,0.4);
\draw[solid,dashed] (-1.1,1.6)--(-1.1,0.4);
\draw[solid,dashed] (0,0.4) ellipse (1.1 and 0.495);
% dimensions from original: hole diameter 20, edge 20, height 10
\node[font=\normalsize,fill=white,inner sep=1.5pt] at (0,1.72) {20};
\node[font=\normalsize] at (0,-0.9) {20};
\node[font=\normalsize,right=4pt] at (2.2,1.0) {10};
\end{tikzpicture}
